% =====================================================================
\section{Discussion and outlook}\label{sec:discussion}
% =====================================================================

\subsection{The rate-independent limit and its regime of validity}
\label{sec:ri_validity}

Rate independence is a modelling idealisation, and we adopt it
deliberately rather than as an established fact. It is the
\emph{quasi-static limit} of epigenetic evolution---the regime in which
the chromatin state either equilibrates quickly relative to the
variation of its micro-environment, or remains effectively frozen except
when a threshold is crossed---and it plays here the role that perfect
plasticity or the play operator play in mechanics: questionable in
general, yet capturing the essential phenomenology (threshold
activation, memory, and hysteresis) while yielding a robust,
thermodynamically consistent and predictive skeleton. The price of the
idealisation is explicit. The pure rate-independent model has \emph{no}
intrinsic time scale, so it cannot represent rate-dependent relaxation,
and it predicts the \emph{absence of ratcheting}: repeating an identical
loading cycle leaves the same residual state each time, and the maximal
damage and residual scar depend only on the extreme values of the
loading, not on its rhythm. This is a sharp, falsifiable prediction: an
experimentally observed dependence of the epigenetic scar on the
\emph{frequency} of an intermittent stimulus at fixed amplitude would
signal a genuine viscous (rate-dependent) contribution, quantified by
the parameter~$\varepsilon$ of the viscous extension discussed in
\cref{sec:beyond_ri}.

\subsection{Relation to the quasi-potential landscape paradigm}
\label{sec:quasipotential}

The dominant mathematical realisation of Waddington's epigenetic
landscape~\cite{Waddington1957,Huang2012} constructs it as a
\emph{quasi-potential} from the stationary distribution of a stochastic
gene-regulatory dynamics~\cite{Wang2008,Ferrell2012}. A central lesson
of that programme is that gene-network dynamics is generically a
\emph{non-equilibrium, non-gradient} system: the deterministic drift
does not derive from a single potential, and a non-vanishing rotational
probability flux, sustained by energy input and responsible for
thermodynamic dissipation, is an essential ingredient. This is in
apparent tension with the present framework, which postulates a stored
energy~$E$ and a gradient-type (energetic) evolution. We do not claim
that the full gene-regulatory dynamics is a gradient flow. Rather, RIE
describes the slow, hysteretic \emph{epigenetic-mark} subsystem in the
quasi-static limit, where an effective energetic description is natural
and where the memory of the loading history, not the fast rotational
flux, dominates the response. In this sense RIE is complementary to the
quasi-potential picture: it trades the kinetic and flux content of the
stochastic landscape for an exact thermodynamic bookkeeping (the first
and second laws hold by construction), analytical thresholds, and the
robustness of the variational structure. Clarifying, for a given
biological system, when the quasi-static energetic reduction is
legitimate---and how the effective energy relates to the underlying
quasi-potential---is an important open question.

\subsection{Relation to kinetic bistable-switch models}
\label{sec:kinetic_models}

A large body of work models epigenetic and cell-fate decisions with
kinetic bistable switches: systems of ordinary differential equations
with cooperative feedback whose bifurcation structure produces
bistability and hysteresis~\cite{Huang2007,Ferrell2012}. In the specific
case that motivates our companion paper, coupled reversible and
irreversible switches have been used to explain the
epithelial--mesenchymal transition and its hysteresis~\cite{Tian2013,
CeliaTerrassa2018}. The hysteresis produced by such models is
\emph{kinetic}---the loop depends on the rate at which the control
parameter is swept---whereas the hysteresis of RIE is
\emph{rate-independent}: the loop is a geometric object whose enclosed
area equals the dissipation per cycle, independent of the rhythm of the
stimulus (\cref{sec:hysteresis}). The two viewpoints are complementary.
Kinetic models resolve the molecular circuitry and its time scales at
the cost of many rate constants and of evolution laws that are, in
practice, prescribed by hand; RIE fixes instead the \emph{structure} of
the evolution from three thermodynamic ingredients, obtaining
closed-form thresholds and unconditional consistency, at the cost of
suppressing the kinetics.

\subsection{Beyond rate independence}\label{sec:beyond_ri}

The framework is the rate-independent member of a broader family, and
the route out of it is already present in the paper. Relaxing the
$1$-homogeneity of the dissipation potential---for instance by adding a
quadratic term $\tfrac{\varepsilon}{2}\|\dot{\bs q}\|^2$---turns the
Biot inclusion into a rate-dependent (viscous), doubly-nonlinear
gradient flow. This is exactly the viscous regularisation used to
\emph{construct} the balanced-viscosity solutions of \cref{sec:bv}: the
rate-independent model is the $\varepsilon\to0$ skeleton of a
one-parameter family, and a finite $\varepsilon$ reintroduces relaxation,
frequency dependence and ratcheting~\cite{mielke2012bv,MielkeRoubicek}.
A systematic study of the finite-$\varepsilon$ regime, and the
identification of $\varepsilon$ from frequency-resolved experiments, is a
natural next step. Two further extensions are of direct biological
relevance: \emph{vectorial} state variables ($n\ge2$), needed to
represent several coupled epigenetic marks, for which the transition
path at a switch becomes genuinely path-dependent and the
balanced-viscosity selection is essential; and \emph{non-symmetric}
dissipation potentials, $\Psi(-\dot q)\ne\Psi(\dot q)$, which model the
direction-dependent resistance of methylation versus demethylation and
the associated asymmetry of the forward and reverse thresholds.

\subsection{Interpretation, identifiability and the thermodynamic
statements}\label{sec:interpretation}

Two caveats delimit the reading of the framework. First, the
ingredients $W$, $\bs\ell$ and $\rho$ are \emph{effective},
phenomenological objects, not directly measured energies; their value is
that they are, in principle, identifiable from macroscopic observables,
since the closed-form thresholds and residual scar
(\cref{sec:experiments}, and, for the bistable model, the companion
paper) express them in terms of quantities accessible to a suitably
designed loading--unloading experiment. Second, the first- and
second-law statements of \cref{sec:thermo} are exact \emph{within} the
thermodynamics of internal variables of the model; they should be read
as an effective, structural thermodynamics---with $E$ a Helmholtz-type
effective free energy and $\mathcal D$ an effective entropy
production---rather than as a literal accounting of molecular energetics.
Their role is to guarantee that any RIE model, whatever the choice of
landscape and dissipation, is internally consistent with the balance
principles, which is precisely the robustness that motivates the
framework.
